
\SGARefTarget{sga7.sga7i.eviii.n2-2}{2.2}. Consider now, in addition to $P$ and $Q$, a commutative group
$G$.  In (\SGARefLink{sga7.sga7i.eviii.n1-1-3}{1.1.3}) we defined a canonical homomorphism, functorial in
$P,Q$, and $G$,
\begin{equation}\tag{\SGARefTarget{sga7.sga7i.eviii.n2-2-1}{2.2.1}}
 \mathrm{Biext}^1(P,Q;G)
 \longrightarrow
 \operatorname{Hom}\bigl(\operatorname{Tor}_1(P,Q),G\bigr).
\end{equation}
On the other hand, if $n$ is a positive integer (respectively, if
$\ell$ is a prime number), the homomorphism $\alpha^n_{P,Q}$ of
(\SGARefLink{sga7.sga7i.eviii.n2-1-1}{2.1.1}) (respectively, the homomorphism $\alpha^{(\ell)}_{P,Q}$ of
(\SGARefLink{sga7.sga7i.eviii.n2-1-13-2}{2.1.13.2})) gives a homomorphism
\[
 \operatorname{Hom}\bigl(\operatorname{Tor}_1(P,Q),G\bigr)
 \longrightarrow
 \operatorname{Hom}({}_nP\otimes{}_nQ,G)
\]
(respectively, a homomorphism
\[
 \operatorname{Hom}\bigl(\operatorname{Tor}_1(P,Q),G\bigr)
 \longrightarrow
 \operatorname{Hom}\bigl(
   T_\ell(P)\otimes T_\ell(Q),T_\ell(G)
 \bigr)).
\]
These last homomorphisms are also functorial in $P,Q,G$.
Composing them with (\SGARefLink{sga7.sga7i.eviii.n2-2-1}{2.2.1}), we obtain canonical homomorphisms,
functorial in $P,Q,G$,
\begin{equation}\tag{\SGARefTarget{sga7.sga7i.eviii.n2-2-2}{2.2.2}}
 \varphi^n=\varphi^n_{P,Q}:
 \mathrm{Biext}^1(P,Q;G)
 \longrightarrow
 \operatorname{Hom}({}_nP\otimes{}_nQ,G),
\end{equation}
\begin{equation}\tag{\SGARefTarget{sga7.sga7i.eviii.n2-2-3}{2.2.3}}
 \varphi^{(\ell)}=\varphi^{(\ell)}_{P,Q}:
 \mathrm{Biext}^1(P,Q;G)
 \longrightarrow
 \operatorname{Hom}\bigl(
   T_\ell(P)\otimes T_\ell(Q),T_\ell(G)
 \bigr).
\end{equation}
For
\begin{equation}\tag{\SGARefTarget{sga7.sga7i.eviii.n2-2-4}{2.2.4}}
 \xi\in\mathrm{Biext}^1(P,Q;G),
\end{equation}
denote its image under (\SGARefLink{sga7.sga7i.eviii.n2-2-2}{2.2.2}) (respectively, under (\SGARefLink{sga7.sga7i.eviii.n2-2-3}{2.2.3})) by
$\varphi^n_\xi$ (respectively, by $\varphi^{(\ell)}_\xi$):
\begin{equation}\tag{\SGARefTarget{sga7.sga7i.eviii.n2-2-5}{2.2.5}}
 \varphi^n_\xi:{}_nP\otimes{}_nQ
 \longrightarrow G
 \qquad\text{(or, equivalently, }{}_nG\text{)},
\end{equation}
\begin{equation}\tag{\SGARefTarget{sga7.sga7i.eviii.n2-2-6}{2.2.6}}
 \varphi^{(\ell)}_\xi:
 T_\ell(P)\otimes T_\ell(Q)
 \longrightarrow T_\ell(G).
\end{equation}

\SGARefTarget{sga7.sga7i.eviii.n2-2-7}{2.2.7}. Plainly, knowledge of $\varphi^{(\ell)}_\xi$ is equivalent
to knowledge of the $\varphi^n_\xi$ for all $n=\ell^\nu$, with
$\nu\geq0$.  For fixed $n$, knowledge of $\varphi^n_\xi$ is
equivalent to knowledge, on the image of
\[
 \alpha^n_{P,Q}:{}_nP\otimes{}_nQ\longrightarrow\uTor_1(P,Q),
\]
of the homomorphism
\begin{equation}\tag{\SGARefTarget{sga7.sga7i.eviii.n2-2-8}{2.2.8}}
 \uTor_1(P,Q)\longrightarrow G
\end{equation}
which is the image of $\xi$ under (\SGARefLink{sga7.sga7i.eviii.n2-2-1}{2.2.1}).

Observe that if $P=nP$ and $Q=nQ$, then the preceding image is
$\,{}_n\uTor_1(P,Q)$.  This follows readily by considering the exact
sequence (\SGARefLink{sga7.sga7i.eviii.n1-3-5}{1.3.5}) and the analogous exact sequence for $Q$.
Furthermore, if $P$ and $Q$ are divisible, that is, if $nP=P$ and
$nQ=Q$ for every positive integer $n$, one readily checks that
$\uTor_1(P,Q)$ is torsion, that is,
\[
 \uTor_1(P,Q)=\varinjlim_n{}_{n}\uTor_1(P,Q).
\]
Thus, in this case, knowledge of the $\varphi^n_\xi$ for all $n$ is
equivalent to knowledge of (\SGARefLink{sga7.sga7i.eviii.n2-2-8}{2.2.8}), whereas knowledge of
$\varphi^{(\ell)}_\xi$ for a given prime $\ell$ is equivalent to
knowledge of the restriction of (\SGARefLink{sga7.sga7i.eviii.n2-2-8}{2.2.8}) to the $\ell$-primary
component of $\uTor_1(P,Q)$.

We now make explicit the symmetry properties of the pairings (\SGARefLink{sga7.sga7i.eviii.n2-2-5}{2.2.5})
and (\SGARefLink{sga7.sga7i.eviii.n2-2-6}{2.2.6}) associated with biextensions.  First observe that there is
a commutative diagram
\begin{equation}\tag{\SGARefTarget{sga7.sga7i.eviii.n2-2-9}{2.2.9}}
\begin{tikzcd}[column sep=large,row sep=huge]
 \mathrm{Biext}^1(P,Q;G)
   \arrow[r,"\sim"]
   \arrow[d]&
 \mathrm{Biext}^1(Q,P;G)
   \arrow[d]\\
 \operatorname{Hom}\bigl(\uTor_1(P,Q),G\bigr)
   \arrow[r,"\sim"]&
 \operatorname{Hom}\bigl(\uTor_1(Q,P),G\bigr),
\end{tikzcd}
\end{equation}
where the horizontal arrows are the symmetry arrows (the first of
which was defined in VII \SGARefLink{sga7.sga7i.evii.n2-7}{2.7}), and the vertical arrows are (\SGARefLink{sga7.sga7i.eviii.n2-2-1}{2.2.1}).
The commutativity follows from that of the two constituent squares in
the fuller diagram
\begin{equation}\tag{\SGARefTarget{sga7.sga7i.eviii.n2-2-10}{2.2.10}}
\begin{tikzcd}[column sep=large,row sep=large]
 \mathrm{Biext}^1(P,Q;G)
   \arrow[r,"\sim"]
   \arrow[d]&
 \mathrm{Biext}^1(Q,P;G)
   \arrow[d]\\
 \operatorname{Ext}^1\!\left(
   P\mathbin{\overset{\mathbb L}{\otimes}}Q,G
 \right)
   \arrow[r,"\sim"]
   \arrow[d]&
 \operatorname{Ext}^1\!\left(
   Q\mathbin{\overset{\mathbb L}{\otimes}}P,G
 \right)
   \arrow[d]\\
 \operatorname{Hom}\bigl(\uTor_1(P,Q),G\bigr)
   \arrow[r,"\sim"]&
 \operatorname{Hom}\bigl(\uTor_1(Q,P),G\bigr).
\end{tikzcd}
\end{equation}
The first vertical arrow in each column is the canonical isomorphism
of VII \SGARefLink{sga7.sga7i.evii.n3-6-5}{3.6.5}, and the second is the evident homomorphism.  The second
square in (\SGARefLink{sga7.sga7i.eviii.n2-2-10}{2.2.10}) is trivially commutative: by definition, the
symmetry isomorphism
\[
 P\mathbin{\overset{\mathbb L}{\otimes}}Q
 \xrightarrow{\ \sim\ }
 Q\mathbin{\overset{\mathbb L}{\otimes}}P
\]
induces the symmetry isomorphism
$\uTor_1(P,Q)\simeq\uTor_1(Q,P)$, and, for a variable complex
$L_\bullet$, the canonical isomorphism
\[
 \operatorname{Ext}^1(L_\bullet,G)
 \simeq
 \operatorname{Hom}\bigl(H_1(L_\bullet),G\bigr)
\]
is functorial in $L_\bullet$.

As for the commutativity of the first square of (\SGARefLink{sga7.sga7i.eviii.n2-2-10}{2.2.10}), it follows
by making explicit simultaneously the symmetry isomorphisms and the
vertical isomorphisms defined in VII \SGARefLink{sga7.sga7i.evii.n3-6-5}{3.6.5}.  We leave the details of
the verification to the reader.  One uses the fact that the symmetry
isomorphism
\[
 L_\bullet(P)\otimes L_\bullet(Q)
 \simeq
 L_\bullet(Q)\otimes L_\bullet(P)
\]
induces in degree $1$ the isomorphism
\[
 L_0(P)\otimes L_1(Q)+L_1(P)\otimes L_0(Q)
 \simeq
 L_0(Q)\otimes L_1(P)+L_1(Q)\otimes L_0(P),
\]
which, on corresponding summands, is given by the ordinary symmetry
isomorphism of the underlying ungraded modules (and hence introduces
no sign).

In view of the commutativity of (\SGARefLink{sga7.sga7i.eviii.n2-2-9}{2.2.9}), the anticommutativity of
(\SGARefLink{sga7.sga7i.eviii.n2-1-10-1}{2.1.10.1}), in terms of the $\varphi^n$ or
$\varphi^{(\ell)}$, becomes the following relation.

\medskip
\noindent\underline{Proposition \SGARefTarget{sga7.sga7i.eviii.n2-2-11}{2.2.11}.}
\underline{The following diagrams are anticommutative:}
\begin{equation}\tag{\SGARefTarget{sga7.sga7i.eviii.n2-2-11-1}{2.2.11.1}}
\begin{tikzcd}[row sep=large,column sep=huge]
 \mathrm{Biext}^1(P,Q;G)
   \arrow[rr,"\sim"]
   \arrow[dd,"{\varphi^n_{P,Q}}"']&&
 \mathrm{Biext}^1(Q,P;G)
   \arrow[dd,"{\varphi^n_{Q,P}}"]\\
 &- &\\
 \operatorname{Hom}({}_nP\otimes{}_nQ,G)
   \arrow[rr,"\sim"]&&
 \operatorname{Hom}({}_nQ\otimes{}_nP,G),
\end{tikzcd}
\end{equation}
\begin{equation}\tag{\SGARefTarget{sga7.sga7i.eviii.n2-2-11-2}{2.2.11.2}}
\begin{tikzcd}[row sep=large,column sep=huge]
 \mathrm{Biext}^1(P,Q;G)
   \arrow[rr,"\sim"]
   \arrow[dd,"{\varphi^{(\ell)}_{P,Q}}"']&&
 \mathrm{Biext}^1(Q,P;G)
   \arrow[dd,"{\varphi^{(\ell)}_{Q,P}}"]\\
 &- &\\
 \operatorname{Hom}\bigl(
   T_\ell(P)\otimes T_\ell(Q),T_\ell(G)
 \bigr)
   \arrow[rr,"\sim"]&&
 \operatorname{Hom}\bigl(
   T_\ell(Q)\otimes T_\ell(P),T_\ell(G)
 \bigr),
\end{tikzcd}
\end{equation}
\underline{where the horizontal arrows are the symmetry arrows.}

We record the important corollary.

\medskip
\noindent\underline{Corollary \SGARefTarget{sga7.sga7i.eviii.n2-2-12}{2.2.12}.}
\underline{Suppose that $P=Q$, and suppose that $\xi$ is the class of
a symmetric (respectively, antisymmetric) biextension of $(P,P)$ by
$G$; that is, suppose ${}^s\xi=\xi$ (respectively,
${}^s\xi=-\xi$), where ${}^s\xi$ denotes the image of $\xi$ under
symmetry.  Then the bilinear maps $\varphi^n_\xi$, on
${}_nP\times{}_nP$, and $\varphi^{(\ell)}_\xi$, on
$T_\ell(P)\times T_\ell(P)$, are antisymmetric (respectively,
symmetric).}

\medskip
\noindent\underline{Remark \SGARefTarget{sga7.sga7i.eviii.n2-2-13}{2.2.13}.}
Suppose that $\xi$ is symmetric.  If $n$ (respectively, $\ell$) is
odd, then \SGARefLink{sga7.sga7i.eviii.n2-2-12}{2.2.12} shows that $\varphi^n_\xi$ (respectively,
$\varphi^{(\ell)}_\xi$) is in fact an alternating form.  This
conclusion is not valid in general without the restriction that $n$
(respectively, $\ell$) be odd, as is already seen when $P$ is the
ordinary group $\mathbb Z/2\mathbb Z$ and
$G=\mathbb Z/2\mathbb Z$.

Suppose, however, that $2P=P$.  This is the case in particular when
$P$ is an abelian scheme, working as usual in the fppf or fpqc
topology.  The preceding forms are then still alternating.  It is
enough to prove that $\varphi^{2^\nu}_\xi(x,x)=0$.  Since $2P=P$, a
given element $x\in{}_{2^\nu}P(S)$ may be written as $x=2y$, with
$y\in{}_{2^{\nu+1}}P(S)$.  Hence
\[
 \varphi^{2^\nu}_\xi(x,x)
 =2\varphi^{2^{\nu+1}}_\xi(y,y)=0,
\]
because $\varphi^{2^{\nu+1}}_\xi$ is antisymmetric.
